\section{Experiments}
\label{sec:experiments}

We report results from a pilot evaluation designed to establish initial baselines, validate the benchmark design, and investigate the example-versus-counterexample performance gap.

\subsection{Evaluation Setup}

\paragraph{Model.}
We evaluate Claude Sonnet~4.6~\citep{anthropic2024claude}, a frontier language model from Anthropic, using the Claude Code CLI (\texttt{claude -p}) in print mode with default settings.
We focus on a single strong model to enable deep qualitative analysis rather than superficial comparison across many models; multi-model evaluation is planned for future work.

\paragraph{Subset selection.}
From the full 332-problem benchmark, we select a stratified random subset of 100 problems.
The stratification ensures proportional representation across chapter, difficulty level, and question type, with minimum constraints of 20 counterexample problems, 20 Level~1 problems, and 10 Level~3 problems (when available) to provide statistical power for key comparisons.
The resulting subset contains 70 examples, 23 counterexamples, 4 computations, and 3 proofs, distributed as: 47 Basic Concepts, 21 Computations, 17 Moduli, 15 Cohomology; and 22 Level~1, 67 Level~2, 11 Level~3.

\paragraph{Prompt.}
Each problem is presented using the prompt template described in Section~\ref{sec:benchmark} (full prompt in Appendix~\ref{app:prompt}).

\paragraph{Scoring.}
Responses were scored using a two-stage pipeline.
First, an automated LLM-based scorer evaluated each response against the known solution using the 0/0.5/1 scale.
The scorer was presented with the problem, known solution, and model response, and asked to classify the response as correct (1.0), partially correct (0.5), or incorrect (0.0), along with a failure mode classification (F1--F6, defined in Section~\ref{sec:analysis}).
Problems where the model timed out (no response within 600 seconds) were automatically scored as 0.0.
We discuss the limitations of this scoring approach in Section~\ref{sec:conclusion}.

\subsection{Pilot Results}

\paragraph{Overall performance.}
Table~\ref{tab:pilot-results} presents the pilot results by difficulty level.

\begin{table}[h]
\centering
\caption{Pilot evaluation results: Claude Sonnet~4.6 on 100-problem stratified subset. Accuracy is the proportion of scores $= 1.0$; Partial includes scores $= 0.5$. Eighteen problems timed out (counted as Incorrect).}
\label{tab:pilot-results}
\begin{tabular}{lccccc}
\toprule
& \textbf{$n$} & \textbf{Correct} & \textbf{Partial} & \textbf{Incorrect} & \textbf{Accuracy} \\
\midrule
Overall & 100 & 78 & 4 & 18 & 0.78 \\
\midrule
Level 1 (Routine) & 22 & 21 & 0 & 1 & 0.95 \\
Level 2 (Non-trivial) & 67 & 54 & 1 & 12 & 0.81 \\
Level 3 (Deep) & 11 & 3 & 3 & 5 & 0.27 \\
\bottomrule
\end{tabular}
\end{table}

The overall accuracy of 78\% masks a steep difficulty gradient: 95\% at Level~1, 81\% at Level~2, and 27\% at Level~3.
The 68-percentage-point drop from Level~1 to Level~3 demonstrates that \ecagbench{} effectively discriminates between routine graduate-level constructions and deep problems requiring synthesis of multiple results.
Including partial credit, the model achieves an average score of 80\%.

\paragraph{Timeout analysis.}
Of the 100 problems, 18 (18\%) timed out without producing a response within 600 seconds.
The timeout rate increases steeply with difficulty: 4.5\% at Level~1, 17.9\% at Level~2, and 45.5\% at Level~3.
Among the 82 problems that received responses, the accuracy is 95.1\% (78/82), suggesting that when the model can respond, it tends to produce correct constructions---but harder problems more frequently cause the model to fail silently.

\paragraph{Performance by question type.}
Table~\ref{tab:pilot-by-type} reveals a pronounced gap between examples and counterexamples.

\begin{table}[h]
\centering
\caption{Pilot results by question type, showing the example--counterexample gap.}
\label{tab:pilot-by-type}
\begin{tabular}{lccccc}
\toprule
\textbf{Question Type} & \textbf{$n$} & \textbf{Correct} & \textbf{Partial} & \textbf{Incorrect} & \textbf{Accuracy} \\
\midrule
Example & 70 & 56 & 1 & 13 & 0.80 \\
Counterexample & 23 & 15 & 3 & 5 & 0.65 \\
Computation & 4 & 4 & 0 & 0 & 1.00 \\
Proof & 3 & 3 & 0 & 0 & 1.00 \\
\bottomrule
\end{tabular}
\end{table}

The model achieves 80\% accuracy on examples but only 65\% on counterexamples---a gap of 15 percentage points.
This gap is consistent with our hypothesis that counterexample construction requires deeper understanding of theorem boundaries (Section~\ref{sec:dual-structure}).
Among the 4 partial-credit responses, 3 were on counterexample problems, further suggesting that counterexamples elicit more fragile reasoning.

\paragraph{Performance by chapter.}
Table~\ref{tab:pilot-by-chapter} shows results across the four represented chapters.

\begin{table}[h]
\centering
\caption{Pilot results by chapter.}
\label{tab:pilot-by-chapter}
\begin{tabular}{lcccc}
\toprule
\textbf{Chapter} & \textbf{$n$} & \textbf{Correct} & \textbf{Accuracy} & \textbf{Partial rate} \\
\midrule
Cohomology & 15 & 14 & 0.93 & 0.00 \\
Basic Concepts & 47 & 37 & 0.79 & 0.04 \\
Computations & 21 & 16 & 0.76 & 0.00 \\
Moduli & 17 & 11 & 0.65 & 0.12 \\
\bottomrule
\end{tabular}
\end{table}

Cohomology has the highest accuracy (93\%), possibly reflecting the structured and well-documented nature of cohomological computations.
Moduli \& Deformation has the lowest accuracy (65\%) and the highest partial-credit rate (12\%), consistent with the more abstract and less standardized nature of the constructions required.

\subsection{Case Studies}

We present three case studies illustrating representative model behaviors.

\paragraph{Case Study 1: Successful counterexample construction (\ecagid{0043}).}
This Level~2 problem asks for the precise difference between ``locally of finite type'' and ``of finite type'' for morphisms of schemes, along with a separating example.
The model correctly identified quasi-compactness as the distinguishing property and constructed $X = \coprod_{n=0}^{\infty} \mathbb{A}^n_k \to \Spec(k)$ as a morphism that is locally of finite type but not of finite type.
The verification correctly showed that each affine piece is a finitely generated $k$-algebra but the source is not quasi-compact.
This demonstrates the model's strength on well-established counterexamples in the graduate curriculum.

\paragraph{Case Study 2: Property confusion on counterexample (\ecagid{0343}, F2).}
This Level~3 problem concerns a $\mathbb{Z}/2\mathbb{Z}$-action on $\mathbb{P}^3$ and asks to verify geometric properties of the action and compute an invariant ring.
The model substituted a different group action ($g \cdot [x:y:z:w] = [-x:-y:z:w]$) from the one specified in the problem setup, then verified properties for this substituted action.
While the algebraic calculations for the substituted action were correct, the mismatch with the intended action constitutes a property confusion error: the construction is valid but answers a different question.

\paragraph{Case Study 3: Timeout on deep problem (\ecagid{0363}).}
This Level~3 problem asks to compute genus-0 Gromov--Witten invariants using Schubert calculus on Grassmannians.
The model timed out after 600 seconds without producing any response.
This pattern---where the model fails silently on problems requiring deep synthesis of advanced techniques (Gromov--Witten theory, Schubert calculus, quantum cohomology)---accounted for 45.5\% of Level~3 problems, making timeout the dominant failure mode at the highest difficulty level.

\subsection{Comparison with Existing Benchmarks}

Direct comparison with other benchmarks is difficult due to differences in domain and evaluation methodology.
However, we note several reference points:
\begin{itemize}[nosep]
  \item FrontierMath~\citep{glazer2024frontiermath} reported that frontier models solve fewer than 2\% of their research-level problems, though their problems are designed to be at the frontier of mathematical knowledge.
  \item On the MATH benchmark~\citep{hendrycks2021measuring}, frontier models now achieve 90+\% accuracy, indicating that competition-level mathematics is largely solved.
  \item On PutnamBench~\citep{tsoukalas2024putnambench}, formal theorem proving success rates remain low for the hardest problems.
\end{itemize}

\ecagbench{} occupies an intermediate position: the 78\% overall accuracy and 27\% Level~3 accuracy suggest that routine constructive algebraic geometry is accessible to frontier models, but problems requiring deep synthesis of multiple advanced techniques remain challenging.
The steep difficulty gradient (95\% $\to$ 81\% $\to$ 27\%) indicates that \ecagbench{} effectively spans this transition.
